\section{Introduction}
\label{sec:intro}

Converting document images into faithful, machine-readable structure --- body text in reading order, formulas as \LaTeX{}, tables as HTML --- underpins retrieval-augmented generation, digital archiving, accessibility, and the construction of training corpora for foundation models. The current state of the art is held by specialized vision--language models (VLMs) of 0.9B--241B parameters~\cite{paddleocrvl,mineru25,glmocr}, and by GPU pipeline systems such as MinerU~\cite{mineru} and Marker~\cite{marker} that chain multi-hundred-megabyte detectors and recognizers, typically including a VLM component for OCR or formulas. Both families presuppose datacenter or at least discrete-GPU deployment.

A large class of real deployments cannot meet that presupposition: on-premise processing of confidential documents, batch digitization on commodity office hardware, edge and offline settings, and cost-sensitive high-volume ingestion. For these, the relevant question is not \emph{what is the best score achievable}, but \emph{how much of that score survives a hard budget} --- no GPU, a few gigabytes of RAM, seconds per page.

This paper answers that question with PRISM, a document parsing pipeline designed under three simultaneous constraints: (i)~CPU-only execution, (ii)~total model weights under 250\,MB, and (iii)~single-digit seconds per page. PRISM contains no generative model larger than 20M parameters. Its accuracy instead comes from an unusually careful orchestration layer around four small specialists --- a layout detector (RT-DETR, 124\,MB ONNX), a text OCR stack ($\sim$11\,MB), a distilled formula recognizer (Texo, 20M parameters), and a table-structure model (SLANet-plus, 7.4\,MB) --- with deterministic geometric reasoning doing the work that end-to-end systems leave to attention: formula band-splitting by ink profiles, fraction-bar fusion, dense-grid ``answer key'' detection, recursive XY-cut reading order, spanning-cell resolution, and photometric/geometric normalization of camera-captured pages.

On OmniDocBench v1.6~\cite{omnidocbench} --- the de facto standard benchmark, 1651 pages across nine real-world document types in English and Chinese, scored by the official public harness --- PRISM reaches an overall score of 78.5\todo{update to v14}, ahead of the widely deployed Marker pipeline (78.4) and within 8 points of the MinerU GPU pipeline (86.5), while every system above it uses at least an order of magnitude more compute. In a controlled same-machine CPU comparison, PRISM dominates the sub-300M ``efficient VLM'' class (SmolDocling~\cite{smoldocling}, GraniteDocling) on every accuracy metric while running $8$--$14\times$ faster in $0.6\times$ the memory.

Our contributions are:
\begin{itemize}
  \item \textbf{A complete, reproducible recipe} for CPU-only document parsing at $<$250\,MB, including all routing, thresholding, and post-processing logic, validated end-to-end on the official OmniDocBench harness.
  \item \textbf{Deterministic structure recovery algorithms} that substitute for model capacity: ink-mass formula grouping (display-formula CDM $56.9 \rightarrow 79.0$ with zero new parameters), recursive XY-cut reading order for multi-column and dense pages, and token-interpolation table assembly with spanning-cell emission.
  \item \textbf{A measurement-first methodology}: we show that a large fraction of apparent recognition failure on this benchmark is attributable to \emph{evaluation-pipeline interactions} (matcher pairing, delimiter nesting, marginalia penalties, un-compilable macro emissions), quantify each, and fix them at the pipeline level; we also report every negative result (detector ensembling, frequency-domain moir\'e removal, threshold sweeps) with numbers.
  \item \textbf{An efficiency frontier study} placing pipelines and small VLMs on the same axes --- accuracy, latency, peak RAM, and model size --- on identical hardware.
\end{itemize}
